\documentclass[10pt,a4paper]{amsart}
\usepackage[margin=19mm]{geometry}
\usepackage{amsmath,amssymb,mathtools}
\usepackage[scr=rsfs,cal=euler]{mathalfa}
\usepackage{xfrac}
\usepackage[unicode,hidelinks,bookmarks=false]{hyperref}
\newcommand{\ie}{i.e.\ }
\newcommand{\cf}{cf.\ }
\newcommand{\resp}{respectively\ }
\newcommand{\Subsection}{\subsection}
\providecommand{\nobreakdash}{\nobreak\hbox{-}}
\setlength{\emergencystretch}{2em}
\multlinegap0pt
\usepackage{bm}
\usepackage{bbm}
\usepackage{dsfont}
\usepackage{xspace}
\usepackage{cancel}
\usepackage{mathtools}
\usepackage{amsthm}
\makeatletter
\@ifundefined{theorem}{\newtheorem{theorem}{Theorem}}{}
\@ifundefined{proposition}{\newtheorem{proposition}[theorem]{Proposition}}{}
\makeatother
\title[Toric periods for a $p$-adic quaternion algebra]{Toric periods for a $p$-adic quaternion algebra}
\author{U. K. Anandavardhanan, Basudev Pattanayak}
\date{Source: arXiv:2304.09765v1 (CC BY 4.0)}
\begin{document}
\maketitle

\noindent\emph{Source and licence.} Toric periods for a $p$-adic quaternion algebra. Version arXiv:2304.09765v1 (CC BY 4.0), distributed under the Creative Commons Attribution 4.0 International licence, as recorded in the arXiv licence metadata for this exact version and shown on the abstract page https://arxiv.org/abs/2304.09765v1. Full rights notice: licenses/arxiv-2304.09765-CC-BY-4.0.txt.\par\smallskip

\noindent\textit{Editorial scope.} Counted unit: the Proposition whose source label is \texttt{prop2} in the article (source0.tex lines 313--317, source label \texttt{prop2}), delivered together with the article's own complete original proof of it (source0.tex lines 319--397, one \texttt{proof} environment of 79 lines). No mathematics is written, repaired or re-derived: statement and proof are the article's own text line for line. Required context retained uncounted: none. Every definition, display and label of those lines is retained unchanged. Omitted from this package and not redistributed: 1--312, 318--318, 398--942. Nothing inside a retained excerpt is omitted or altered: every retained body is delivered exactly as the article's lines are written, so no omission receipt of any kind accompanies this package. Citations: the retained excerpts carry no citation command at all, so no bibliography entry of the article is transcribed and no reference is redistributed. Reference adaptation: none was needed and none was made; every reference inside the delivered unit resolves inside it, so no unresolved reference or citation is produced. Printed slips kept literal: no printed word, symbol, hypothesis or number of the counted statement or of its proof was altered. Layout and class: the article's own class is \texttt{amsart} with packages including amscd, amssymb, latexsym, amsmath, xy, anysize, mathpazo, color; the wrapper is the lane's pinned \texttt{amsart} editorial wrapper at 10pt on A4, which restates the theorem-like environments the retained text needs and adds the article's own macro definitions that the retained text needs (none), with extra wrapper packages (bm, bbm, dsfont, xspace, cancel, mathtools). The delivered numbering is the unit's own. Assets: no retained span contains a figure environment, a graphics-inclusion command, a PDF-page inclusion, a table or a TikZ picture; no image file is copied into this package, the package declares no asset dependency and no image file is redistributed with it. Licence: version arXiv:2304.09765v1 of this article is distributed under the Creative Commons Attribution 4.0 International licence, as recorded in the arXiv licence metadata for this exact version and shown on the abstract page https://arxiv.org/abs/2304.09765v1. A full rights notice is delivered at licenses/arxiv-2304.09765-CC-BY-4.0.txt. Characters: every retained excerpt is pure ASCII, so the delivered engine, XeLaTeX with its default Latin Modern text font, reports no missing character.
\begin{proposition}\label{prop2}
For $1 \leq i \leq \frac{q-1}{2}$, pick the element $g=\zeta^i \in D^\times$ and consider a smooth character $\kappa_g$ of $K^\times$ such that on $\kappa_g|_{g^{-1}Hg \cap K^\times} = \psi^{{\#}^g}$. Then the character $\kappa_g$ occurs with multiplicity one in the restriction of $\pi_\psi$ to $K^\times$ and $\kappa_g|_{F^\times} = \psi|_{F^\times}$. Let $\ell \in \{0,1\}$ be such that $\kappa_g(\varpi) = (-1)^\ell \psi(\varpi)$. The vector
\[v_{\kappa_g} = \sum\limits_{\underline{\alpha} \in S^{\prime \prime} } \kappa_g(1 + \sum\limits_{r} \alpha_r \varpi^r) \left[ f_{\underline{\alpha}, i}^\prime + (-1)^\ell \psi(\zeta^{(1-i)(q+1)})  f_{\underline{\alpha}, q+1-i}^\prime \right]\] 
lies in the $\kappa_g$-isotypic component of the representation $\pi_\psi$.
\end{proposition}

\begin{proof}
The first assertion is straightforward Mackey theory. We now show that $v_{\kappa_g}$ is a $\kappa_g$-invariant vector. To this end, observe that $v_{\kappa_g}(x) = 0$ unless fix $x = h \cdot d_{\underline{\delta}, i}^\prime$ or 
$x = h \cdot d_{\underline{\delta}, {q+1-i}}^\prime$ for $h \in H$, $\underline{\delta}=(\delta_1,\dots,\delta_{\frac{m-1}{2}}) \in \mathcal S^{\prime\prime}$. Clearly,
\[
v_{\kappa_g}(x) = 
\begin{cases}
\psi^\#(h) \cdot \kappa_g(1 + \sum\limits_{r} \delta_r \varpi^r)  &\text{if $x = h \cdot d_{\underline{\delta}, i}^\prime $,} \\
(-1)^\ell \psi(\zeta^{(1-i)(q+1)}) \psi^\#(h) \cdot \kappa_g(1 + \sum\limits_{r} \delta_r \varpi^r)  &\text{if $x = h \cdot d_{\underline{\delta}, q+1-i}^\prime $,} \\
0 &\text{otherwise.}
\end{cases}
\]
We write $z \in K^\times$ as 
\[z= a \cdot \varpi^{j_0} \cdot (1 + \sum\limits_{\substack{j \geq 1  \\    j \text{ is odd} }}\beta_j \cdot  \varpi^j )\]
for some $a \in F^\times$, $j_0 \in \{0,1\}$, and $\beta_j \in \mathbb F$ for all $r$. Clearly, when $a \in F^\times \subset K^\times $, we have 
\[\pi_\psi(a) v_{\kappa_g}(x)=v_{\kappa_g}(xa)=\psi(a) \cdot v_{\kappa_g}(x) = \kappa_g(a) v_{\kappa_g}(x).\]
Now we check the action of the uniformizer $\varpi \in K^\times$ and elements of the form 
\[z^\prime = 1 + \sum\limits_{\substack{j \geq 1  \\    j \text{ is odd} }}\beta_j \cdot  \varpi^j \]
with $\beta_r \in \mathbb F$ for all $r$. Without loss of generality we may take $x =  h \cdot d_{\underline{\delta}, i}^\prime$, as the proof in the case $x =  h \cdot d_{\underline{\delta}, q+1-i}^\prime$ being similar. Then, 
\begin{align*}
\pi_\psi(\varpi) v_{\kappa_g}(x) &= v_{\kappa_g}(h \cdot d_{\underline{\delta}, i}^\prime  \cdot \varpi) \\
&= v_{\kappa_g}(h \cdot \varpi \cdot \zeta^{(i-1)(q+1)} \cdot d_{\underline{\delta}, q+1-i}^\prime) \\
&= \psi^\#(h) \cdot \psi(\varpi) \cdot \psi(\zeta^{(i-1)(q+1)}) \cdot \kappa_g(1 + \sum\limits_{r} \delta_r \varpi^r)  (-1)^\ell \psi(\zeta^{(1-i)(q+1)}) \\
&= (-1)^\ell\psi(\varpi)  \cdot  \psi^\#(h) \kappa_g(1 + \sum\limits_{r} \delta_r \varpi^r)  \\
&= \kappa_g(\varpi) \cdot v_{\kappa_g}(x).
\end{align*}
Now we check the action of $z^\prime$. We have 
\begin{align*}
\pi_\psi(z^\prime) v_{\kappa_g}(x) &= v_{\kappa_g}(h \cdot d_{\underline{\delta}, i}^\prime \cdot (1 + \sum\limits_{\substack{j \geq 1  \\    j \text{ is odd} }}\beta_j \cdot \varpi^j )) \\
&= v_{\kappa_g}\left( h (1 + \sum\limits_{\substack{ r = 1  \\  r \text{ is odd} }}^{\frac{m-1}{2}} \delta_r \zeta^{i(1-q)} \varpi^r)  \cdot (1 + \sum\limits_{\substack{j \geq 1  \\    j \text{ is odd} }}\beta_j  \zeta^{i(1-q)} \varpi^j ) \zeta^i \right).
\end{align*}
Let $m^*$ be the largest odd integer in $\{1,2,...,\frac{m-1}{2}\}$. Then the element 
\[\textbf{y} = (1 + \sum\limits_{\substack{r = 1  \\    r \text{ is odd} }}^{m^*} \delta_r \zeta^{i(1-q)} \varpi^r) (1 + \sum\limits_{\substack{j \geq 1  \\    j \text{ is odd} }}\beta_j  \zeta^{i(1-q)} \varpi^j ) \zeta^i\]
can be written as 
\[\textbf{y } = b_0 \cdot h_0 \cdot d^\prime_{\underline{\delta}^{\prime},i}\] 
where,
\[b_0 = 1 + \sum\limits_{\substack{r = 2 \\    r \text{ is even} }}^{m^*+1} \xi_r  \varpi^r \]
with 
\[\xi_r = \sum\limits_{\substack{t=1  \\    t \text{ is odd} }}^{r-1} \delta_{r-t}\beta_t \]
for each even $r \leq m^*+1$, 
and
\[\underline{\delta}^{\prime} = (\delta_1^{\prime},...,\delta_{\frac{m-1}{2}}^{\prime}) \in S^{\prime\prime}\]
with 
\[\delta_r^{\prime} = \delta_r + \beta_r  - \sum\limits_{\substack{t =1 \\    t \text{ is odd} }}^{r-2} \delta_t^{\prime} \xi_{r-t},\]
for each odd $r \leq m^*$, 
and
\[h_0 = 1 + \sum\limits_{\substack{j =m^*+2  \\    j \text{ is odd} }}^{m} \epsilon_j \zeta^{i(1-q)}\varpi^j + \sum\limits_{\substack{k =m^* +3 \\    k \text{ is even} }}^{m-1} \xi_k \varpi^k +  \cdots \]
such that
\[\epsilon_j = \beta_j -\sum\limits_{\substack{t=1  \\    t \text{ is odd} }}^{m^*} \delta_t^\prime \xi_{j-t}- \sum\limits_{\substack{t=m^* +2  \\    t \text{ is odd} }}^{j-2} \epsilon_t \xi_{j-t}, \]
for odd $j \geq \frac{m+1}{2}$,
and
\[\xi_k = \sum\limits_{\substack{t=1  \\    t \text{ is odd } }}^{m^*}\delta_t \beta_{k-t} -\sum\limits_{\substack{t=m^*+3  \\    t \text{ is even} }}^{k-2} \xi_t \xi_{k-t} - \sum\limits_{\substack{t=m^*+2  \\    t \text{ is odd} }}^{k-1} \delta_{k-t}^\prime \epsilon_{t},\]
for even $k \geq \frac{m+1}{2}$.
Denote 
\[\tilde{h} = g^{-1}h_0 g = 1 + \sum\limits_{\substack{j =m^*+2  \\    j \text{ is odd} }}^{m} \epsilon_j \varpi^j + \sum\limits_{\substack{k =m^* +3 \\    k \text{ is even} }}^{m-1} \xi_k \varpi^k +  \cdots.\]
Note that $\tilde{h} \in g^{-1}Hg \cap K^\times$. We have
\begin{align*}
 \pi_\psi(z^\prime) v_{\kappa_g}(x) &= v_{\kappa_g}(h \cdot b_0 \cdot h_0 \cdot d_{\underline{\delta}^\prime,i}^\prime) \\
&= \psi^\#(h) \cdot \psi(b_0) \cdot \psi^{\#^g}(\tilde{h}) \cdot \kappa_g(1 + \sum\limits_{\substack{r = 1  \\    r \text{ is odd} }}^{m^*} \delta_r^{\prime} \varpi^r) ,\\
&= \psi^\#(h) \cdot \kappa_g \left(b_0 \cdot \tilde{h} \cdot (1 + \sum\limits_{\substack{r = 1  \\    r \text{ is odd} }}^{m^*} \delta_r^{\prime} \varpi^r)\right).\\
\end{align*}
Observe that 
\begin{align*}
z^\prime \cdot (1 + \sum\limits_{\substack{r = 1  \\    r \text{ is odd} }}^{m^*} \delta_r \varpi^r) 
&= (1 + \sum\limits_{\substack{j \geq 1  \\    j \text{ is odd} }} \beta_j  \varpi^j ) \cdot (1 + \sum\limits_{\substack{r = 1  \\    r \text{ is odd} }}^{m^*} \delta_r \varpi^r) \\
&= b_0 \cdot (1 + \sum\limits_{\substack{j =m^*+2  \\    j \text{ is odd} }}^{m} \epsilon_j \varpi^j + \sum\limits_{\substack{k =m^* +3 \\    k \text{ is even} }}^{m-1} \xi_k \varpi^k +  \sum\limits_{t \geq m+1} \gamma_t \varpi^t) \cdot (1 + \sum\limits_{\substack{r = 1  \\    r \text{ is odd} }}^{m^*} \delta_r^{\prime} \varpi^r) 
\end{align*}
for some $\gamma_t \in \mathbb F$ for $t \geq m+1$. Since $ \kappa_g|_{(1 + \mathcal{P}_K^{m+1})}=1$, we have 
\begin{align*}
\kappa_g \left(b_0 \cdot \tilde{h} \cdot (1 + \sum\limits_{\substack{r = 1  \\    r \text{ is odd} }}^{m^*} \delta_r^{\prime} \varpi^r)\right) 
&=  \kappa_g \left( b_0 \cdot (1 + \sum\limits_{\substack{j =m^*+2  \\    j \text{ is odd} }}^{m} \epsilon_j \varpi^j + \sum\limits_{\substack{k =m^* +3 \\    k \text{ is even} }}^{m-1} \xi_k \varpi^k)\cdot (1 + \sum\limits_{\substack{r = 1  \\    r \text{ is odd} }}^{m^*} \delta_r^{\prime} \varpi^r)\right) \\
&= \kappa_g \left (z^\prime \cdot (1 + \sum\limits_{\substack{r = 1  \\    r \text{ is odd} }}^{m^*} \delta_r \varpi^r) \right) = \kappa_g(z^\prime)  \kappa_g(1 + \sum\limits_{r} \delta_r \varpi^r). 
\end{align*}
Therefore,
\begin{align*}
 \pi_\psi(z^\prime) v_{\kappa_g}(x) &= \kappa_g(z^\prime) \cdot \psi^\#(h) \kappa_g(1 + \sum\limits_{r} \delta_r \varpi^r) \\
 &= \kappa_g(z^\prime) \cdot v_{\kappa_g}(x).
\end{align*}
Hence, $\pi_\psi(z) \cdot v_{\kappa_g}= \kappa_g(z) \cdot v_{\kappa_g}$ for all $z \in K^\times$.
\end{proof}

\end{document}
